\subsection{伪模}

给定一个含或不含单位元的伪环 A，A 上的左伪模是一个交换群 E（加法记法），它以 A 为算子集，且满足 §1, 第 1 号, 定义 1 的公理 \((\mathbf{M}_{\mathrm{I}})\) , \((\mathbf{M}_{\mathrm{II}})\) 和 \((\mathbf{M}_{\mathrm{III}})\) 。A 上的右伪模类似地定义。

设 \(A'\) 是给 A 添加单位元所得的环。若 E 是 A 上的一个左伪模，则在 E 上关联一个左 \(A'\) -模结构，其方式是对所有 \(x \in E\) 及每个元素 \((n, a) \in A'\) ，写下

\[
(n, a). x = n x + a x.\tag{2}
\]

公理 \((\mathbf{M}_{\mathrm{I}})\) 到 \((\mathbf{M}_{\mathrm{IV}})\) （§1, 第 1 号, 定义 1 的）立即得到验证；此外，把这个模结构的算子集限制到 \(\{0\} \times A\) （与 A 等同）上，我们就在 E 上重新得到最初给定的伪模结构。

对于 E 的子集 M，要使 M 成为伪模 E 的带算子子群（此时诱导结构显然也是 A 上的左伪模结构），当且仅当 M 是相伴 A'-模 E 的子模，且此子 A'-模与伪模 M 相伴。进而，商 A'-模 E/M 与带算子商群 E/M 相伴，而后者显然是 A 上的伪模。

若 E、F 是 A 上的两个伪模，则带算子群同态 \(E \rightarrow F\) 与分别相伴于伪模 E 和 F 的 A'-模之间的 A'-线性映射 \(E \rightarrow F\) 完全相同。若 \((\mathbf{E}_{t})_{t \in I}\) 是 A 上的一族伪模，则带算子群 \(\prod_{i\in I}E\) 和 \(\bigoplus_{i\in I}E_{i}\) 是 A 上的伪模，且相伴的 A'-模分别是相伴 A'-模 \(E_{i}\) 的积与直和。对伪模的逆极限与直极限也有类似的结果。因此，A 上伪模的理论可以完全归结为 A'-模的理论。

